\uuid{TtGI}
\titre{Équation droite/plan}
\auteur{bodin ; hanani ; mzari}
\organisation{exo7}
\licence{CC BY-NC-SA 4.0}
\datecreate{2020-01-20}
\niveau{L1}
\module{Géométrie}
\chapitre{Géométrie de l'espace}
\sousChapitre{Plans}
\difficulte{1}
\theme{Plans, équation droite, plan}
\exo7qcmid{lille-1.geometrie-espace.004}

\contenu{
\texte{
Pour ces questions, l'espace est muni d'un repère orthonormé direct  $(O, \vec{i}, \vec{j}, \vec{k})$.
}

\begin{question}
\qclassification{141.03}
\qtags{motcle=équation droite/plan}
Soit $P$ le plan passant par $A(1,1,0)$ et  de vecteur normal $\vec{n}(1,-1,1)$. Quelles sont les assertions vraies ?
\begin{answers}
    \bad{Une équation cartésienne de $P$ est $x-y+z=1$.}

    \good{Une équation cartésienne de $P$ est $x-y+z=0$.}

     \good {Une représentation paramétrique de $P$ est :
 $$\left\{\begin{array}{ccl}x&=&t-s\\y&=&t\\ z&=&s,\quad (t,s\in \Rr) \end{array}\right.$$}

    \good{Une représentation paramétrique de $P$ est :
$$\left\{\begin{array}{ccl}x&=&t\\y&=&s\\ z&=&s-t,\quad (t,s \in \Rr)\end{array}\right.$$}

\end{answers}
\begin{explanations}
Une équation de $P$ est de la forme : $x-y+z+a=0$ et on cherche $a$ pour que $A$ appartienne à $P$. On résout cette équation pour trouver une représentation paramétrique.
\end{explanations}
\end{question}
}
